\section{State the teaching subset before writing a parser}
The exercise uses one image, \texttt{P6}, a Maxval token of exactly \texttt{255}, and
width and height from one through sixty-four. These are local bounds chosen for a small
framebuffer and a bounded input arena. They are not limits of the general PPM format.
The final Maxval separator is one LF byte in this exercise. The writer emits the
canonical eleven-byte style illustrated above, with decimal dimensions taking as many
digits as necessary.

The reader will accept header spacing and comments in its declared header gaps, but
will not search for comments inside the raster. Width and height use unsigned decimal
tokens without signs or leading zeros. Reject an empty token, a zero dimension, a digit
that would exceed the bound, and a nonseparator after a token. Bound the complete header
to 256 bytes. A comment without a terminating line feed must reach that bound or the
end of input and fail; it must not scan unrelated simulator memory for a terminator.

Require exactly one complete raster and no trailing bytes. This makes a truncated row,
an extra byte and a second image explicit failures for this exercise. A more general
reader could expose a next-image offset; changing that policy requires a new contract
and tests rather than silently accepting extra input.

The display path preserves sample codes directly. It does not implement a color-management
conversion or use an image viewer as the numerical oracle. Test the RGB bytes and display
words first. A viewer is useful as an additional check that the saved file can be opened.
